\label{sec:ExistingImplementation}
Explaining the full \textit{C++} implementation by \underline{\hyperlink{https://git-crysp.uwaterloo.ca/iang/prac.git}{Vadapalli et al.}} \cite{vadapalli2023duoram} exceeds the scope of this research.
Instead, the code will be explained by highlighting the most important fragments and scenarios.
The author's extension can be found \underline{\hyperlink{https://github.com/FelixWeik/PrivateRandomAccessComputations}{here}}.

\subsection*{Preliminaries}
\label{sub:ImplementationPreliminaries}
\subsubsection*{Representing DPFs} \label{subsub:implementationRepresentingDPFs}
Both, Distributed Point Functions (\ref{sub:DistributedPointFunctions}) and Distributed Comparison Functions (\ref{sub:DistributedComparisonFunctions}), stem from the same parent class \code{dpf}.
Again, recall point-function and comparison-function represented as binary trees.
Thus, the most trivial representation of a single point function is saving the seed of the root node and generating the whole tree every time the point function is used.
The base class therefore saves the seed of the root node and implements the functionality of descending from a given seed to one of its children.
Also, it stores the flag bits of the leaf level with bit $i$ being leaf $i$.
This is equivalent to a single vector $t_b$ as explained in section \ref{sec:3pDuoram} and will XOR to $e_i$ when combined with the  the other party's counterpart.
The code further distinguishes Distributed Point Functions for memory accesses, called \textit{RDPF}, and Distributed Comparison Functions labelled as \textit{CDPF}.

RDPFs inherit from the basic DPF implementation.
It expands its parent with the \code{LeafInfo} struct.
This stores the final correction word and the sum of scaling to receive a share, which can be either additive or a XOR share, of the message $M$.
Thus, combined with the flags of the leaf-vector implemented by the parent, the RDPF is stored in an already evaluated form holding both $t_b$ and $v_b$.
%TODO Knoten werden als 128-bit werte gespeichert => warum ist das so => erklären!!!

%TODO das ist aes verschlüsselt -> soll man das miterklären oder ist das vorerst egal weil es ja keine rolle spielt?
The code stores RDPFs as triples, called \code{RDPFTriple}.
As the name suggests, it holds three RDPFs, each with the same target, again, either XOR or additively shared.
To generate a RDPF-triple including three rdpfs on a given target, one uses its constructor which includes the complete functionality to do so.
This construction makes it easy to keep track of the used RDPFs.

CDPFs too, inherit from the given DPF construction.
They are, however, stored as tuples.
Both CDPFs have the same target \code{as\_target} for additive shares, and \code{xs\_target} for XOR shares.
A tuple of CDPFs allows the client to compare a given value to 0 as described in \ref{sub:CDPFConstruction}.

\subsubsection*{types.hpp} \label{subsub:implementationTypeshpp}
This file defines basic macros, aliases, structs, and their operators.

The most prominent type-aliases are:
\begin{itemize}
    \item \code{nbits\_t}: 
    This value counts the number of bits in a value, which means \code{value\_t}.
    Thus, the chosen datatype needs to be large enough to store \textit{VALUE\_BITS} bits.
    The default value is \code{uint8\_t}.
    \item \code{bit\_t}:
    Chooses how a single bit is represented, usually in the \code{RegBS} struct.
    The default datatype for this alias is \code{bool}.
    \item \code{address\_t}: 
    Deceides the size of the addresses in the secret-shared memory.
    Is set to \code{uint32\_t}.
    \item \code{DPFnode}: 
    Implements the datatype of a DPFnode.
    It is of type \code{\_\_m128i}.
    Leafs are also implemented using this type.
    Evaluation on leaf nodes, however, splits the vector in a lower and upper half which makes it two 64-bit vectors.
    \item \code{value\_t}: 
    Deceides the datatype of values in the secret-shared memory.
    The default value is \code{uint64\_t}.
\end{itemize}
The most prominent macros are:
\begin{itemize}
    \item {VALUE\_BITS}: 
    This will deceide the size of \code{value\_t}. 
    The default value is 64.
    \item {ADDRESS\_MAX\_BITS}: 
    Deceides the size of \code{address\_t}.
    The default value is 32.
\end{itemize}
The most prominent structs are:
\begin{itemize}
    \item \code{RegAS}: 
    Represents a register holding an additive share of type \code{value\_t}.
    \item \code{RegXS}: 
    Represents a register holding an XOR share of type \code{value\_t}.
    \item \code{RegBS}: 
    Represents a bit-share of type \code{bit\_t} which is a type-alias for a bool.
\end{itemize}
\subsubsection*{Intel AES-NI} \label{subsub:implementationAesni}
\subsubsection*{duoram} \label{subsub:implementationDuoram}
\subsubsection*{mpcio} \label{subsub:implementationMpcio}
\subsubsection*{Communication} \label{subsub:implementationCommunication}

\subsection*{Scenarios} \label{sub:implementationScenarios}
\label{sub:ImplementationScenarios}
\subsubsection*{Preprocessing} \label{subsub:implementationPreprocessing}
%TODO wie wird das in files gespeichert und wo findet man die?
\subsubsection*{Cyclic Shifting} \label{subsub:implementationCyclicShifting}
\subsubsection*{Online Phase} \label{subsub:implementationOnlinePhase}
%TODO wie werden die DPFs (RDPF, CDPF) aufgerufen
\subsubsection*{Reading} \label{subsub:implementationReading}
\subsubsection*{Writing} \label{subsub:implementationWriting}